%%%%%%%%%%%%%%%%%%%%%%%%%%  teacher_draft_updated.tex  %%%%%%%%%%%%%%%%%%%%%%%%%%
%
% Teacher's draft + (a) simplified Section 5 (Theory),
%                    (b) proofread Section 6 (Experiments),
%                    (c) minor factual fixes in Setup.
%
% Compile: pdflatex -> bibtex -> pdflatex -> pdflatex
%%

\documentclass[twoside,leqno,twocolumn]{article}
\usepackage[letterpaper]{geometry}
\usepackage{siamproceedings}

\usepackage[T1]{fontenc}
\usepackage{amsfonts}
\usepackage{graphicx}
\usepackage{epstopdf}
\usepackage{enumitem}
\usepackage{algorithmic}
\ifpdf
  \DeclareGraphicsExtensions{.eps,.pdf,.png,.jpg}
\else
  \DeclareGraphicsExtensions{.eps}
\fi

\newsiamremark{remark}{Remark}
\newsiamremark{hypothesis}{Hypothesis}
\crefname{hypothesis}{Hypothesis}{Hypotheses}
\newsiamthm{claim}{Claim}

\usepackage{amsopn}
\DeclareMathOperator{\diag}{diag}

\usepackage{amsmath,amssymb}
\usepackage{amsthm}
\usepackage{algorithm}
\usepackage{algpseudocode}
\usepackage{booktabs}
\usepackage{graphicx}
\usepackage{subcaption}
\usepackage[hidelinks]{hyperref}

\graphicspath{
  {./}
  {./results/}
  {./results/multidata/}
  {./OverclusterDRG_Paper/}
  {./OverclusteringRKC/}
  {./overclustering_me/}
}

\newcommand{\OPT}{\mathrm{OPT}}
\newcommand{\rz}{r_z}
\newcommand{\CL}{C_L}
\newcommand{\Copt}{C^{*}_{Q}}
\newcommand{\rkc}{R$k$C}

\begin{document}

\newcommand\relatedversion{}

\title{\Large OverclusterDRG: a Deterministic, LP-Free Heuristic for
Robust $k$-Center\relatedversion}

\date{}
\maketitle

\fancyfoot[R]{\scriptsize{Copyright \textcopyright\ 20XX by SIAM\\
Unauthorized reproduction of this article is prohibited}}

% ----------------------------------------------------------------
\begin{abstract}
The Robust $k$-Center problem asks for $k$ centers minimising the
maximum distance to all but $z$ outlier points. Practical methods sit
at two extremes: the Charikar et al.\ greedy $3$-approximation delivers
quality at quadratic cost, while the RKC-Gonzalez heuristic is fast
but has no guarantee; Ding et al.\ lie between them, but their trial
count grows linearly in~$z$. We present OverclusterDRG, a
deterministic, LP-free algorithm: run Gonzalez for $k+O$ steps to
build a candidate pool, then select $k$ centers by greedy
initialisation and $1$-swap local search. On 36 configurations of UCI
Adult ($N$ up to $20{,}000$, exact LP bounds), OC-Local attains a mean
LP ratio of $1.147$, beating Charikar in $33/36$ configurations at a
median $7\times$ speedup and Ding in $35/36$ at $8\times$ speedup; all
paired tests are significant at $p<10^{-5}$. The algorithm rests on an
exact phase transition: the pool contains a $3\cdot\OPT$ solution
whenever $O\ge z$, and the ratio is unbounded for $O<z$, so $O=z$ is
the smallest valid overcount.
\end{abstract}
% ----------------------------------------------------------------


\section{Introduction}

The $k$-center problem asks for $k$ centers minimising the
maximum distance from any point to its nearest center. Because it
optimises a worst-case distance, it is the natural model for
applications with strict service guarantees---emergency facility
placement, sensor coverage, network latency bounds. Gonzalez
\cite{gonzalez1985} gave an $O(Nk)$-time greedy algorithm with a
tight $2$-approximation: start from any point, and at each step add
the point farthest from the current centers.

One outlier breaks everything. A distant anomaly forces the algorithm
to place an entire center there, leaving a true cluster unserved.
Robust $k$-Center (\rkc) fixes this by allowing up to $z$ points to
be discarded; the radius is then evaluated on the remaining $N-z$
points. The robust formulation handles noise well but is harder to
solve: choosing which $z$ points to discard is itself a combinatorial
problem layered on top of NP-hard clustering.

The gap between what the theory provides and what practice needs is
large. Existing algorithms with guarantees are either too slow or
randomised:
\begin{itemize}[leftmargin=*]
  \item \textbf{Charikar et al.}~\cite{charikar2001}: a $3$-approximation
  built on optimal-radius guessing and disk intersection. It works, but
  scales as $O(N^{2.12})$ empirically---hundreds of seconds on
  moderately large data.
  \item \textbf{Exact LP}: the LP relaxation gives the tightest known
  lower bound but is out of reach beyond $N\approx2{,}000$; at
  $N=20{,}000$ a single solve takes up to 94 hours.
  \item \textbf{Ding et al.}~\cite{ding2019}: a $2$-approximation via
  randomised greedy, needing $T=\lceil(z{+}1)\ln100\rceil$ trials for
  99\% confidence. At $z=20\%$ of $N=20{,}000$ that is $18{,}423$
  trials, taking 431 seconds. Two runs give different answers.
  \item \textbf{Naive Gonzalez}: run $k$ farthest-point steps and
  discard the $z$ farthest assigned points. Fast at $O(Nk)$ but has no
  approximation guarantee; one outlier can misdirect a center. Our
  experiments show LP ratios of $1.28$--$1.92$ depending on dataset.
\end{itemize}

No algorithm is simultaneously deterministic, LP-free, fast at large
$N$ and $z$, and tied to $\OPT$.

We present \textbf{OverclusterDRG} --- \textbf{Over}clustering
followed by a \textbf{D}eterministic \textbf{R}obust \textbf{G}reedy
selection. The key idea: Gonzalez wastes steps on outliers, so give it
$O$ spare steps. With $O=z$ extra steps the pool absorbs all outliers
and still holds a representative for each true cluster. We select $k$
centers from this pool by greedy initialisation followed by $1$-swap
local search, producing a single deterministic pass.

\smallskip\noindent\textbf{Contributions.}
\begin{enumerate}[leftmargin=*]
  \item \textbf{Algorithm}: OverclusterDRG, deterministic and LP-free,
  running in $O(Nkz)$ per sweep, completing the hardest configuration
  ($N=20{,}000$, $z=20\%$) in 57 seconds.
  \item \textbf{Empirical quality}: mean LP ratio $1.147$ on UCI Adult
  (exact LP), beating Charikar in $33/36$ cases ($7\times$ faster) and
  Ding in $35/36$ ($8\times$ faster). The ordering holds across
  Diabetes ($1.155$) and Covertype ($1.171$). On 400 small instances,
  $1$-swap matches the exhaustive pool optimum in $97.2\%$ of cases.
  \item \textbf{Theory}: an exact phase transition at $O=z$---the pool
  contains a $3\cdot\OPT$ subset for $O\ge z$ and no finite guarantee
  holds below it. One extra Gonzalez step gives a free certified lower
  bound $\OPT\ge D_{k+z+1}/2$ at every~$N$.
\end{enumerate}


% ----------------------------------------------------------------
\section{Problem Setup}
\label{sec:setup}

Let $(X,d)$ be a finite metric space with $|X|=N$. Given integers
$k\ge1$ and $z\ge0$, find $C\subseteq X$ with $|C|=k$ and
$S\subseteq X$ with $|S|\le z$ minimising
\[
  R(C,S) = \max_{x\in X\setminus S} d(x,C), \quad
  d(x,C)=\min_{c\in C}d(x,c).
\]
The optimum is $\OPT=\min_{|C|=k,\,|S|\le z}R(C,S)$.

For fixed $C$ the optimal $S$ is just the $z$ farthest points, so we
write $\rz(C,X)$ for the $(z{+}1)$-th largest value in
$\{d(x,C):x\in X\}$ and $\OPT=\min_{|C|=k}\rz(C,X)$.

\paragraph{LP lower bound.}
For a fixed radius $r$, define $N(j,r)=\{i \mid d(i,j)\le r\}$ as the
set of points close enough to serve as a center for $j$.
Introduce variables: $y_i\in[0,1]$ (center opened at $i$);
$x_{ij}\in[0,1]$ (point $j$ assigned to center $i$);
$c_j\in[0,1]$ (point $j$ is covered, i.e.\ not an outlier).
\begin{align*}
\sum_{i\in N(j,r)} x_{ij} &\ge c_j \quad\forall j
  &&\text{($j$ covered only if assigned within radius $r$)}\\
x_{ij} &\le y_i \quad\forall i,j
  &&\text{(assign only to open centers)}\\
\sum_i y_i &\le k
  &&\text{(at most $k$ centers)}\\
\sum_j c_j &\ge N-z
  &&\text{(at most $z$ outliers)}\\
x_{ij},\,y_i,\,c_j &\in [0,1]
\end{align*}
We binary-search over all pairwise distances as candidates for $r$
and solve this feasibility problem at each probe. The smallest
feasible $r$ is $R_{\mathrm{LP}}$, which satisfies
$R_\mathrm{LP}\le\OPT$ and serves as the lower bound for all
reported ratios. The full LP is in Appendix~\ref{app:lp}.

\paragraph{Two optima.}
Let $\OPT_{m,0}=\min_{|C|=m}\max_{x\in X}d(x,C)$ be the standard
$m$-center optimum covering all $N$ points. The relation between
$\OPT_{k+O,0}$ and $\OPT_{k,z}$ is the core of the theory
(Section~\ref{sec:theory}).


% ----------------------------------------------------------------
\section{Related Work}

Gonzalez \cite{gonzalez1985} gave the $2$-approximation for $k$-center
in $O(Nk)$; its \rkc\ adaptation evaluates $\rz$ on the result but
has no guarantee.

Charikar et al.\ \cite{charikar2001} gave the first constant-factor
approximation for outlier clustering; the \rkc\ variant certifies
$3r^{*}$ at quadratic cost. Ding et al.\ \cite{ding2019} achieve a
$2$-approximation via randomised greedy, requiring
$T=\lceil(z{+}1)\ln100\rceil$ trials for 99\% confidence---a
multiplier that grows with $z$.

Local search gives the best bounds for $k$-median: Arya et al.\
\cite{arya2004} proved $(3+2/p)$ for $p$-swap via telescoping sums.
The argument works because $k$-median is a sum of distances; for
$k$-center (a maximum) it breaks, which is why our locality-gap
conjecture (Section~\ref{sec:conjecture}) is open. Goyal and Jaiswal
\cite{goyal2021} give tight FPT approximations for constrained
$k$-center; these target a different (exponential-in-$k$) regime.


% ----------------------------------------------------------------
\section{The OverclusterDRG Algorithm}
\label{sec:algorithm}

\paragraph{Phase 1: pool construction.}
Run Gonzalez for $k+O$ steps. The starting point is the data point
nearest the empirical centroid $\bar{X}$; each subsequent point is
$q_{i+1}=\arg\max_x\min_{j\le i}d(x,q_j)$. Ties are broken
lexicographically by coordinates, making the output independent of
input order. The pool is $Q=\{q_1,\dots,q_{k+O}\}$, built in
$O(N(k{+}O))$. The default overcount is $O=z$---Section~\ref{sec:theory}
proves this is the minimum valid choice.

\paragraph{Phase 2a: forward greedy (OC-Forward).}
Pre-compute the $N\times(k{+}O)$ distance matrix once. Starting from
$\{q_1\}$, repeatedly add the pool point whose inclusion minimises
$\rz$, breaking ties by lowest pool index. Cost: $O(NkO)$.

\paragraph{Phase 2b: $1$-swap local search (OC-Local).}
For each pair $(c_i,q)\in C\times(Q\setminus C)$, scanned in
fixed pool-index order, accept the first swap that strictly decreases
$\rz$; repeat until no improving swap exists. The objective strictly
decreases at each accepted swap and the pool has finitely many
$k$-subsets, so the loop terminates. In practice $5$--$10$ accepted
swaps suffice on every configuration tested.

\begin{algorithm}[t]
\caption{Gonzalez$(X, k)$}
\label{alg:gonzalez}
\begin{algorithmic}[1]
\STATE Start from an arbitrary point. Set $C = \{c_1\}$.
\STATE Add the point farthest from all current centers in $C$.
\STATE Repeat Step 2 until $k$ centers are selected.
\STATE \textbf{return} $C$
\end{algorithmic}
\end{algorithm}

\begin{algorithm}[t]
\caption{OverclusterDRG$(X,k,z,O{=}z)$}
\label{alg:main}
\begin{algorithmic}[1]
\STATE $Q \gets \textsc{Gonzalez}(X,\,k+O)$
  \COMMENT{Phase 1: pool of $k{+}O$ candidates}
\STATE $C_0 \gets \textsc{OC-Forward}(X,Q,k,z)$
  \COMMENT{Phase 2a: greedy initialisation}
\STATE $C \gets \textsc{OC-Local}(X,Q,C_0,z)$
  \COMMENT{Phase 2b: $1$-swap improvement}
\STATE \textbf{return} $C$
\end{algorithmic}
\end{algorithm}

\paragraph{Complexity.}
Phase~2 costs $O(NkO)$ per sweep; with $O=z$ this is $O(Nkz)$ per
sweep, or $O(N^2k)$ when $z=\Theta(N)$. The number of sweeps is
finite but has no polynomial worst-case bound; empirically $5$--$10$
suffice. The dominant memory is the $N\times(k{+}O)$ distance matrix
($O(Nz)$ doubles at $O=z$, about $7.9$\,GB at $N=99{,}000$, $k=20$,
$z=10\%$).


% ----------------------------------------------------------------
\section{Theoretical Analysis}
\label{sec:theory}

The guarantee uses four results: a classical pool bound, a new
comparison of the two optima, the phase transition itself (both
directions), and a free certified lower bound. Two further results
concern what local search finds, and whether $O=z$ is the best
operating point.

\subsection{The phase transition at $O=z$}

\begin{lemma}[Pool guarantee]
\label{lem:pool}
Let $Q=\textsc{Gonzalez}(X,k{+}O)$. There exists a $k$-subset
$C_0\subseteq Q$ with $\rz(C_0)\le\OPT_{k,z}+2\cdot\OPT_{k+O,0}$.
\end{lemma}
\begin{proof}
Let $(C^*,S^*)$ be an optimal \rkc\ solution. For each $c^*_j\in
C^*$ pick the nearest pool point $q_j$; by the Gonzalez
$2$-approximation $d(c^*_j,q_j)\le2\cdot\OPT_{k+O,0}$. Set
$C_0=\{q_1,\dots,q_k\}$ (pad on collisions). For any inlier $x\notin
S^*$ with nearest optimal center $c^*_j$,
\[
  d(x,C_0)\le d(x,c^*_j)+d(c^*_j,q_j)\le\OPT+2\cdot\OPT_{k+O,0}.
\]
Discarding $S^*$ gives $\rz(C_0)\le\OPT+2\cdot\OPT_{k+O,0}$.
\end{proof}

\begin{lemma}[OPT comparison]
\label{lem:optcompare}
For $O\ge z$: $\OPT_{k+O,0}\le\OPT_{k,z}$.
\end{lemma}
\begin{proof}
Take the $k$ optimal \rkc\ centers and add one center at each of the
$z$ outliers; this covers all $N$ points within radius $\OPT$, so
$\OPT_{k+z,0}\le\OPT$. More centers can only help, so
$\OPT_{k+O,0}\le\OPT_{k+z,0}\le\OPT$ for $O\ge z$.
\end{proof}

\begin{theorem}[Sufficiency]
\label{thm:sufficiency}
For $O\ge z$, the pool contains a $k$-subset with $\rz\le3\cdot\OPT$.
\end{theorem}
\begin{proof}
Lemma~\ref{lem:pool}: $\rz(C_0)\le\OPT+2\cdot\OPT_{k+O,0}$.
Lemma~\ref{lem:optcompare}: $\OPT_{k+O,0}\le\OPT$. Therefore
$\rz(C_0)\le\OPT+2\cdot\OPT=3\cdot\OPT$.
\end{proof}

\begin{theorem}[Necessity]
\label{thm:necessity}
For any $k\ge2$, $z\ge1$, $M>0$, there is an instance with
$\OPT=\varepsilon$ on which the pool with $O=z-1$ satisfies
$\min_{C\subseteq Q,|C|=k}\rz(C)\ge M\varepsilon$.
\end{theorem}
\begin{proof}[Proof sketch]
Place $k$ tight clusters of $m>z$ inliers and $z$ outliers at distance
$R\gg\varepsilon$ in orthogonal directions. The outliers are the
farthest points, so Gonzalez picks all $z$ before any second cluster
contributes a representative. With $O=z-1$ the pool has the starting
inlier, the $z$ outliers, and $k-2$ cluster representatives, leaving
one cluster unrepresented. Every $k$-subset either wastes a slot on an
outlier or leaves the orphaned cluster uncovered; since that cluster
has $m>z$ inliers, none of the $z$ budget can discard them all.
Radius $\ge D-2\varepsilon=2M\varepsilon$ follows from $D=2(M{+}1)\varepsilon$.
\end{proof}

\begin{corollary}[Phase transition]
\label{cor:phase}
$O=z$ is the exact threshold: below it the ratio is unbounded ($k\ge2$);
at and above it the pool always contains a $3\cdot\OPT$ subset.
\end{corollary}

The mechanism: Gonzalez wastes exactly one step per outlier. With
$k+O$ total steps and $z$ wasted, the useful steps are $k+O-z$. At
$O=z$ this equals $k$---one representative per cluster, none to spare.

\subsection{A free certified lower bound}

\begin{lemma}[Certified lower bound]
\label{lem:cert}
If $N\ge k{+}z{+}1$, then $\OPT\ge D_{k+z+1}/2$, where $D_t$ is the
selection distance of the $t$-th Gonzalez point.
\end{lemma}
\begin{proof}
The $k+z+1$ Gonzalez points are pairwise at distance $\ge D_{k+z+1}$:
each $q_t$ was chosen at distance $D_t\ge D_{k+z+1}$ from its
predecessors. The optimum discards at most $z$ of them, so at least
$k+1$ survive as inliers. Two of those share an optimal center by
pigeonhole, giving $2\cdot\OPT\ge D_{k+z+1}$.
\end{proof}

Running one extra Gonzalez step yields $D_{k+z+1}$, and every run
therefore carries the certificate $\rz(C)/\OPT\le2\rz(C)/D_{k+z+1}$
with no LP. The certified ratio averages $3.6$ on Adult---valid but
conservative; Section~\ref{sec:expsetup} explains why.

\subsection{Past the threshold: anytime guarantee}
\label{sec:anytime}

Write $f(O)=\min_{C\subseteq Q,|C|=k}\rz(C)$ for the best pool-subset
quality.

\begin{theorem}[Anytime containment]
\label{thm:anytime}
For every $O\ge z$:
\[
  f(O)\;\le\;\bigl(1+2\rho(O)\bigr)\cdot\OPT,
  \qquad\rho(O)=\frac{D_{k+O+1}}{D_{k+z+1}}\;\le\;1.
\]
\end{theorem}
\begin{proof}
Every optimal center $c^*_j$ lies within $D_{k+O+1}$ of some pool
point $q_j$ (by definition of the next Gonzalez selection). Set
$C_0=\{q_1,\dots,q_k\}$. For any inlier $x$ with nearest optimal
center $c^*_j$, $d(x,C_0)\le\OPT+D_{k+O+1}$; discarding $S^*$ gives
$f(O)\le\OPT+D_{k+O+1}\le\OPT+\rho(O)\cdot D_{k+z+1}\le
\OPT+2\rho(O)\cdot\OPT$.
\end{proof}

At $O=z$, $\rho(z)=1$ and the bound recovers $3\cdot\OPT$. As $O$
grows and Gonzalez distances keep falling, $\rho(O)$ decreases and
the bound tightens---post-threshold quality gains are predicted, not
coincidental.

\begin{corollary}[Adaptive pool rule]
\label{cor:adaptive}
Stop Gonzalez at the first $t\ge z$ with $D_{k+t+1}\le\varepsilon
D_{k+z+1}$. The pool then contains a $(1+2\varepsilon)\cdot\OPT$
subset, and the worst case is $t=N-k$ (trivially exact).
\end{corollary}
\begin{proof}
Theorem~\ref{thm:anytime} with $O=t$ and $\rho(t)\le\varepsilon$ by
construction.
\end{proof}

\begin{theorem}[Containment constant $\ge2$]
\label{thm:lowerbound2}
For every $\delta\in(0,1)$ there is a $k=2$ instance on which the pool
with $O=z$ satisfies $f(z)\ge(2-\delta)\cdot\OPT$.
\end{theorem}
\begin{proof}[Proof sketch]
Place cluster $A$ at $\{-1,0,1\}$ and cluster $B$ at $\{b,b+1\}$
with $b=2-\delta$, plus $z$ trap outliers at distance $R\gg 1$.
Gonzalez starts at $+1$, picks all traps, then picks $-1$; cluster
$B$ is never represented. The only feasible $2$-subset has radius
$b=2-\delta$, while $\OPT=1$.
\end{proof}

Combined with Theorem~\ref{thm:sufficiency}, the true worst-case
constant lies in $[2,3]$; closing this interval is open.

\subsection{What local search finds}
\label{sec:localstuff}

\begin{theorem}[$1$-swap has unbounded worst-case gap]
\label{thm:neg}
For any $D>0$ and $O\ge2$ there is an instance where the
greedy-initialised $1$-swap local optimum satisfies
$\rz(\CL)/\OPT>D$.
\end{theorem}
\begin{proof}[Proof sketch]
Take $k$ equal-sized clusters and two singletons at slightly greater
distance. The singletons are the two farthest points, so Gonzalez
places them first in the pool; forward greedy inserts both before
enough cluster representatives, orphaning two clusters. No single swap
covers both, so the local optimum is stuck at radius equal to the
inter-cluster distance $D'$; choosing $D'\gg\varepsilon$ makes the
ratio $D'/\varepsilon$ arbitrarily large.
\end{proof}

This failure requires adversarially equal inter-cluster distances.
On real data it is rare: $1$-swap matches exhaustive search in $97.2\%$
of $400$ test instances (Section~\ref{sec:localvalid}).

\subsection{Locality-gap conjecture}
\label{sec:conjecture}

\begin{conjecture}
On every instance where all optimal clusters have $>z$ points and all
pairwise distances are distinct, every greedy-initialised $1$-swap
local optimum with $O=z$ satisfies
$\rz(\CL)\le(1+2/k)\cdot f(O)$.
\end{conjecture}

The condition is not sufficient to guarantee the bound in general
(the Theorem~\ref{thm:neg} instance satisfies $|V_j|>z$ but still
fails); the obstruction is the flat max-type objective, which lets
two simultaneously uncovered clusters hide behind a radius no single
swap reduces. Pairwise-distinct distances dissolve this exact
flatness. The bound held on all $400$ validation instances
(maximum observed gap $4.82\%$ against $2/k\in\{20\%,40\%\}$).


% ----------------------------------------------------------------
\section{Experiments}
\label{sec:experiments}

\subsection{Setup}
\label{sec:expsetup}

\begin{table}[t]
\caption{Datasets. All features z-score standardised.}
\label{tab:datasets}
\centering
\begin{tabular}{lrrl}
\toprule
Dataset & Rows used & Features & Source \\
\midrule
UCI Adult            & 32{,}561         & 6  & \cite{adult_dataset} \\
Diabetes Indicators  & 99{,}492         & 8  & UCI \\
Covertype            & 30{,}000 of 581k & 10 & \cite{covertype_dataset} \\
HTRU2                & 17{,}898         & 8  & \cite{htru2} \\
\bottomrule
\end{tabular}
\end{table}

Table~\ref{tab:datasets} lists the datasets. The Adult grid uses
$N\in\{500,1000,2000,5000,10000,20000\}$, $k\in\{10,20,50\}$, and
outlier rate $z/N\in\{10\%,20\%\}$, giving 36 configurations.
Diabetes substitutes $N=15{,}000$ for $2{,}000$ (36 configurations);
Covertype uses $N\in\{500,1000,5000,10000,15000\}$ (30
configurations). Five algorithms are compared: RKC-Gonzalez, OC-Forward,
OC-Local, Charikar et al.\ \cite{charikar2001} at its achieved radius,
and Ding et al.\ \cite{ding2019} at $T=\lceil(z{+}1)\ln100\rceil$
trials. We report the ratio of achieved radius to the LP lower bound;
lower is better and $1.000$ is unattainable in general.

\begin{algorithm}[t]
\caption{RKC-Gonzalez$(X, k, z)$}
\label{alg:rkcgonzalez}
\begin{algorithmic}[1]
\STATE Start from an arbitrary point. Set $C = \{c_1\}$.
\STATE Add the $(z{+}1)$-th farthest point from current centers,
       skipping the $z$ most extreme candidates.
\STATE Repeat Step 2 until $k$ centers are selected.
\STATE Drop the $z$ points farthest from their assigned center.
\STATE \textbf{return} $C$ and its radius.
\end{algorithmic}
\end{algorithm}

\paragraph{Preprocessing.}
Adult uses its six numeric features (age, fnlwgt, education-num,
capital-gain, capital-loss, hours-per-week) after z-scoring;
categorical columns are dropped. Covertype is subsampled to
$30{,}000$ rows (NumPy \texttt{default\_rng}, seed $42$), then
z-scored; the same subsample is used for all Covertype configurations.

\paragraph{LP bounds.}
On Adult, the LP is solved exactly at all 36 configurations (Gurobi
12.0.3, feasibility tolerance $10^{-6}$). Runtimes range from
$12$--$1{,}096$\,s at $N\le2{,}000$ to up to $3.4\times10^{5}$\,s (94
hours) at $N=20{,}000$. On Diabetes and Covertype the LP is too
expensive at large $N$; we interpolate $R_\mathrm{LP}$ from the
Adult LP/Gonzalez ratios per $(k,z/N)$ bucket, and flag this in
Table~\ref{tab:multidata}. Because every algorithm within a
configuration shares the same denominator, all within-dataset
comparisons and orderings are exact regardless of whether the LP is
solved or interpolated.

\paragraph{Free certificates.}
Lemma~\ref{lem:cert} gives $\OPT\ge D_{k+z+1}/2$ at no extra cost.
The certified ratio $2\rz(C)/D_{k+z+1}$ averages $3.6$ (Adult),
$3.1$ (Diabetes), $3.3$ (Covertype), and captures $36$--$39\%$ of
the LP bound. The looseness is structural: with $z=\Theta(N)$, the
$(k{+}z{+}1)$-th Gonzalez spacing falls well below $\OPT$ on dense
data. The certificate is tightest at small $N$ and small $z$ (ratio
$\approx2$ at $N=500$) and is best read as a quick sanity bound.

\subsection{UCI Adult: full comparison}

\begin{table}[t]
\caption{LP ratios and runtimes on 36 Adult configurations (exact LP).
Lower is better; $1.000$ is optimal.}
\label{tab:adult}
\centering
\begin{tabular}{lcccccc}
\toprule
Algorithm & Mean & Std & Min & Max & Median\,t & Max\,t \\
\midrule
LP oracle        & 1.000 & ---   & 1.000 & 1.000 & 4{,}586\,s & 338{,}755\,s \\
\textbf{OC-Local}& \textbf{1.147} & 0.047 & 1.088 & 1.265 & 0.36\,s & 57\,s \\
OC-Forward       & 1.158 & 0.047 & 1.104 & 1.283 & 0.26\,s & 52\,s \\
Charikar         & 1.173 & 0.056 & 1.104 & 1.306 & 3.60\,s & 564\,s \\
Ding             & 1.235 & 0.032 & 1.179 & 1.331 & 2.22\,s & 431\,s \\
RKC-Gonzalez     & 1.280 & 0.049 & 1.197 & 1.370 & 0.002\,s & 0.02\,s \\
\bottomrule
\end{tabular}
\end{table}

OC-Local achieves the best mean ratio among all practical algorithms:
$1.147$ against Charikar's $1.173$, Ding's $1.235$, and Gonzalez's
$1.280$. It is also faster: median $0.36$\,s against Charikar's
$3.60$\,s and Ding's $2.22$\,s. Per configuration, OC-Local beats
Charikar in $33/36$ cases at a median $7\times$ speedup, and Ding
in $35/36$ cases at a median $8\times$ speedup. The LP row is shown
for scale---median solve time $4{,}586$ seconds, worst case 94 hours.

\begin{figure}[t]
\centering
\includegraphics[width=0.98\linewidth]{fig3_quality_speed.png}
\caption{Quality--runtime tradeoff on UCI Adult. Each point is one
$(N,k,z/N)$ configuration; lower-left is better. OC-Local (green
diamonds) sits on the quality frontier; no other practical algorithm
matches its radius at comparable runtime.}
\label{fig:qualityspeed}
\end{figure}

\begin{figure}[t]
\centering
\includegraphics[width=0.98\linewidth]{fig1_ratio_vs_N.png}
\caption{LP ratio versus $N$ (averaged over $k$, one panel per outlier
rate). OC-Local stays near $1.13$--$1.16$ throughout; the gap to
Gonzalez and Ding grows with $N$.}
\label{fig:ratiovsN}
\end{figure}

Figure~\ref{fig:qualityspeed} shows the quality--runtime picture;
Figure~\ref{fig:ratiovsN} shows OC-Local's advantage is stable across
scales.

\begin{table}[t]
\caption{Paired tests on per-configuration LP ratios ($n=36$, exact
LP). Positive $\bar{d}$ means OC-Local is better.}
\label{tab:ttests}
\centering
\begin{tabular}{lcccc}
\toprule
Comparison & $\bar d$ & $t$ & $p$ & Wilcoxon $p$ \\
\midrule
vs RKC-Gonzalez & $+0.134$ & $+13.3$ & $<10^{-14}$ & $2.9\!\times\!10^{-11}$ \\
vs Ding         & $+0.089$ & $+14.3$ & $<10^{-15}$ & $5.8\!\times\!10^{-11}$ \\
vs Charikar     & $+0.026$ & $+5.7$  & $<10^{-5}$  & $3.1\!\times\!10^{-8}$ \\
vs OC-Forward   & $+0.012$ & $+7.9$  & $<10^{-8}$  & $2.9\!\times\!10^{-11}$ \\
\bottomrule
\end{tabular}
\end{table}

Table~\ref{tab:ttests} shows none of the gaps are noise. Even the
OC-Local vs OC-Forward comparison---the same algorithm, minus the
local search step---is significant at $p<10^{-8}$: the $1$-swap phase
costs a median $0.10$ seconds and improves the ratio in every one of
the 36 configurations. It should always be applied. SimplifiedDRG
(the pool-free variant) reaches only $1.295$; the pool buys an $11.5\%$
improvement. The 36 configurations share a base dataset, overstating
independence in the paired $t$-tests; the Wilcoxon column (no
distributional assumption) agrees with every conclusion. Restricting
to the $18$ smallest configurations ($N\le2{,}000$, where the LP
solves in seconds) leaves every conclusion intact: OC-Local wins
$17$--$18$ of $18$ with Wilcoxon $p\le1.3\times10^{-4}$.

\paragraph{Runtime scaling.}
Log-log regression against $N$ gives OC-Local $O(N^{1.85})$
($R^2=0.976$) and Charikar $O(N^{2.12})$ ($R^2=0.999$). Fitting
separately within each $(k,z/N)$ stratum leaves OC-Local's exponent
tightly clustered around $1.85$ (none below $1.8$ or above $1.9$,
all $R^2\ge0.95$). Ding's exponent is $O(N^{1.10})$, but the trial
count dominates: at the hardest configuration ($N=20{,}000$, $k=50$,
$z=20\%$) Ding needs $18{,}423$ trials and $431$ seconds while
OC-Local finishes in $57$ seconds.

\subsection{Local search validation}
\label{sec:localvalid}

Theorem~\ref{thm:neg} says $1$-swap can fail badly. How often does it
happen on real data? We ran OC-Local and exhaustive Phase~2 on $400$
independent random subsamples of Adult ($N\le1{,}000$, $k\in\{5,10\}$,
$z\le10$)---the largest regime where exhaustive search is feasible.

\begin{figure}[t]
\centering
\includegraphics[width=0.98\linewidth]{fig5_local_vs_exhaustive.png}
\caption{Gap between $1$-swap and exhaustive search across $400$
instances (left: gap vs pool size; right: histogram of the 11
non-zero gaps). Local search is exact in $97.2\%$ of cases; maximum
gap $4.82\%$.}
\label{fig:localexhaustive}
\end{figure}

OC-Local matched the exhaustive solution in $389$ of $400$ instances
($97.2\%$). The $11$ failures all had the adversarial structure from
Theorem~\ref{thm:neg}: small clusters at near-equal distances. The
maximum gap was $4.82\%$, at pool size $20$ ($k=10$, $z=10$), where
exhaustive search takes $10$ seconds and OC-Local takes $3$\,ms. No
observed gap exceeded $1+2/k$, supporting the conjecture.

\subsection{Pool size: the $O=z$ threshold}
\label{sec:oablation}

Corollary~\ref{cor:phase} predicts a sharp change at $O=z$. We swept
$O$ over $72$ configurations on Diabetes and Covertype, measuring the
best pool-subset radius at each $O$.

\begin{figure}[t]
\centering
\includegraphics[width=0.98\linewidth]{fig_O_ablation.png}
\caption{Radius (normalised to its value at $O=z$) versus $O$ on a
log axis across $72$ configurations. Radius falls steeply while
$O<z$ and flattens near the threshold.}
\label{fig:oablation}
\end{figure}

\begin{figure}[t]
\centering
\begin{subfigure}{0.48\linewidth}
  \includegraphics[width=\linewidth]{fig_phase_transition.png}
  \caption{}
\end{subfigure}
\hfill
\begin{subfigure}{0.48\linewidth}
  \includegraphics[width=\linewidth]{fig_O_vs_zk.png}
  \caption{}
\end{subfigure}
\caption{(a) Radius vs $O$ for three configurations; quality drops
sharply left of $O=z$. (b) Knee $O_\mathrm{knee}/z$ versus $z/k$
($r=-0.60$): more outliers per center push the knee toward $z$.}
\label{fig:phasepair}
\end{figure}

Radius drops steeply below $O=z$ in every configuration
(Figure~\ref{fig:oablation}), as Theorem~\ref{thm:necessity} predicts.
The empirical \emph{knee}---the smallest $O$ within $2\%$ of best pool
quality---sits between $1.0z$ and $1.7z$, mean $1.55z$.

It is worth being clear about what the theory covers and what it does
not. Corollary~\ref{cor:phase} is about what the pool
\emph{contains}---the $3\cdot\OPT$ guarantee switches on at $O=z$.
The knee measures what OC-Local \emph{finds}, which keeps improving
past $z$. This gap between containment and recovery is open.
Theorem~\ref{thm:anytime} gives the theory for the post-threshold
regime: the containment bound is $(1+2\rho(O))\cdot\OPT$, and
$\rho$ falls as Gonzalez distances fall. Measured across the 72
configurations, mean $\rho$ drops from $1.0$ at $O=z$ to $0.76$ at
$O=2z$ and $0.50$ at $O=5z$. Whether the knee coincides with where
$\rho$ flattens is not answered by this data: only $4$ of $72$
configurations show a detectable $\rho$-plateau (correlation $0.69$
at $n=4$, insufficient).

\paragraph{Adaptive rule.}
We ran Corollary~\ref{cor:adaptive} at $\varepsilon\in\{1.0, 0.5,
0.25, 0.1\}$. $\varepsilon=1.0$ stops at $t=z$ by construction
(sanity check passed). Tighter $\varepsilon$ costs real pool: mean
stopping at $3.7z$ ($\varepsilon=0.5$), $6.0z$ ($0.25$), and $6.9z$
($0.1$). Halving the worst-case factor from $3\cdot\OPT$ to
$2\cdot\OPT$ costs roughly $3.7\times$ the pool; $O=z$ remains the
right default.

\begin{table}[t]
\caption{Pool-size settings across $72$ configurations. Gap is relative
to best quality over all tested $O$.}
\label{tab:osettings}
\centering
\begin{tabular}{lcccl}
\toprule
Setting & $\le5\%$ & Mean gap & Max gap & Bound \\
\midrule
$O=z$ & 53/72 & $+3.7\%$ & $+11.8\%$ & $3\cdot\OPT$ \\
$O=2z$ & 69/72 & $+1.3\%$ & $+8.2\%$ & $2.52\cdot\OPT$ \\
Adaptive ($\varepsilon{=}0.5$) & \textbf{72/72} & ${<}1\%$ & ${<}3\%$ & $2\cdot\OPT$ \\
\bottomrule
\end{tabular}
\end{table}

Table~\ref{tab:osettings} shows that $O=z$ lands within $5\%$ of best
in $53/72$ cases; $O=2z$ raises this to $69/72$ at $2\times$ the
runtime; the adaptive rule ($\varepsilon=0.5$) reaches $72/72$ with
the tightest proven bound.

\paragraph{What predicts the knee.}
Across $72$ configurations, $z/k$ is the strongest single predictor of
$O_\mathrm{knee}/z$ ($r=-0.60$, $p<10^{-9}$), with raw $z$
($r=-0.61$) and $N$ ($r=-0.59$) close behind; dimensionality is
uncorrelated ($r=+0.04$, $p=0.72$). Joint OLS ($R^2=0.43$) keeps
$z/k$ as the strongest standardised coefficient ($-0.25$). When each
center covers many outliers, the knee stays near $z$; with few
outliers per center, extra pool helps.

\begin{figure}[t]
\centering
\includegraphics[width=0.98\linewidth]{fig_synthetic.png}
\caption{Synthetic sweeps of $z/k$, cluster separation, and
dimension. The knee stays near $z$ throughout; only $d\ge12$ pushes
it toward $2.5z$.}
\label{fig:synthetic}
\end{figure}

Synthetic data (Figure~\ref{fig:synthetic}) confirm this across two
decades of $z/k$, a range of separations, and dimensions $2$--$20$.

\subsection{Cross-dataset validation}
\label{sec:multidata}

\begin{figure}[t]
\centering
\includegraphics[width=0.98\linewidth]{fig_ratio_both.png}
\caption{LP ratio vs $N$ on Diabetes (left) and Covertype (right).
The algorithm ordering and magnitude replicate the Adult results.}
\label{fig:ratioboth}
\end{figure}

\begin{table}[t]
\caption{LP ratios on secondary datasets. Diabetes and Covertype LP
bounds are interpolated from Adult; see Section~\ref{sec:expsetup}.
Charikar at achieved radius.}
\label{tab:multidata}
\centering
\begin{tabular}{llcccc}
\toprule
Dataset & Algorithm & Mean & Std & Min & Max \\
\midrule
Diabetes    & RKC-Gonzalez & 1.923 & 0.256 & 1.449 & 2.414 \\
($d{=}8$,   & Ding         & 1.227 & 0.048 & 1.143 & 1.330 \\
$n{=}36$)   & Charikar     & 1.201 & 0.044 & 1.145 & 1.313 \\
            & OC-Local     & \textbf{1.155} & 0.043 & 1.093 & 1.275 \\
\midrule
Covertype   & RKC-Gonzalez & 1.671 & 0.173 & 1.418 & 1.960 \\
($d{=}10$,  & Ding         & 1.213 & 0.049 & 1.121 & 1.300 \\
$n{=}30$)   & Charikar     & 1.227 & 0.060 & 1.134 & 1.374 \\
            & OC-Local     & \textbf{1.171} & 0.051 & 1.095 & 1.294 \\
\bottomrule
\end{tabular}
\end{table}

OC-Local has the best mean ratio on both datasets: $1.155$ on
Diabetes and $1.171$ on Covertype, ahead of Charikar ($1.201$ and
$1.227$), Ding ($1.227$ and $1.213$), and Gonzalez ($1.923$ and
$1.671$). The algorithm ordering is the same as on Adult. Gonzalez
degrades more sharply on these datasets, which widens OC-Local's
margin further. On HTRU2 ($N=17{,}000$, $k=10$, $z=10\%$) raw radii
tell the same story: OC-Local $1.79$ in $17$\,s against Gonzalez
$4.21$ and Charikar's achieved radius above OC-Local's in $56$\,s.

\begin{table}[t]
\caption{OC-Local speedup over Ding at full certificate trial count.
OC-Local is faster in all $66$ configurations.}
\label{tab:speedup}
\centering
\begin{tabular}{lcccc}
\toprule
Dataset & Mean & Median & Min & Max \\
\midrule
Diabetes  ($n{=}36$) & $2.0\times$ & $2.0\times$ & $1.3\times$ & $2.4\times$ \\
Covertype ($n{=}30$) & $2.6\times$ & $2.6\times$ & $2.1\times$ & $3.6\times$ \\
\bottomrule
\end{tabular}
\end{table}

OC-Local is faster than Ding in all $66$ configurations by $2$--$3\times$
(Table~\ref{tab:speedup}), smaller than the $7$--$8\times$ on Adult
because each Ding trial is cheap here. The structural advantage is not
speed---it is quality, determinism, and the fact that the trial count
does not grow with $z$.

\subsection{Scalability}
\label{sec:scalability}

\begin{table}[t]
\caption{OC-Local at large $N$ (Diabetes, $k=20$, $z=10\%$).
Values for $N\ge50{,}000$ are projected from the fitted $O(N^{1.85})$
scaling law; Ding times are projected from $O(N^{1.10})$ plus the
trial multiplier.}
\label{tab:scale}
\centering
\begin{tabular}{rccc}
\toprule
$N$ & OC-Local radius & OC-Local time & Ding time \\
\midrule
20{,}000 & 2.060 & 57\,s  & 133\,s \\
50{,}000 & $\sim$2.11 & $\sim$310\,s & $\sim$900\,s \\
99{,}000 & $\sim$2.15 & $\sim$1{,}100\,s & $\sim$3{,}800\,s \\
\bottomrule
\end{tabular}
\end{table}

At $N=99{,}000$ OC-Local is projected to finish in about $18$ minutes;
Ding needs roughly an hour; Charikar's $O(N^{2.12})$ scaling implies
several hours; and the LP is out of reach. OC-Local is the only
algorithm that remains deterministic and carries a formal statement
about what it has searched at this scale.


% ----------------------------------------------------------------
\section{Conclusion}

OverclusterDRG is a fast, deterministic heuristic for Robust
$k$-Center that sits between the quality of Charikar and the speed of
Gonzalez. The key result is the exact phase transition at $O=z$: the
pool contains a $3\cdot\OPT$ subset whenever the overcount equals the
outlier budget, and no finite guarantee holds below it; so $O=z$ is
not a tuning choice but a theoretically pinpointed threshold.

On UCI Adult (exact LP bounds), OC-Local achieves mean LP ratio
$1.147$, beating Charikar in $33/36$ configurations ($7\times$ faster)
and Ding in $35/36$ ($8\times$ faster). The advantage holds across
Diabetes and Covertype and the algorithm remains practical at
$N\approx10^5$ where the LP is unavailable.

Three problems remain open: proving the locality-gap conjecture
(which requires a new technique for max-type objectives); determining
whether the true worst-case constant is $2$ or $3$
(Theorems~\ref{thm:sufficiency} and~\ref{thm:lowerbound2} bound the
interval, but our constructions are repeatedly self-corrected by
farthest-first); and extending the phase-transition analysis to
$k$-median with outliers, where the sum-type objective changes the
geometry entirely.


% ----------------------------------------------------------------
\appendix
\section{Phase-2 Subroutines}
\label{app:pseudocode}

\begin{algorithm}[t]
\caption{OC-Forward$(X, Q, k, z)$}
\label{alg:ocforward}
\begin{algorithmic}[1]
\STATE Set $C = \{q_1\}$, the first pool point.
\STATE For each remaining pool point $q \in Q \setminus C$, compute
       $r_z(C \cup \{q\})$ --- the robust radius if $q$ is added.
\STATE Add to $C$ the point $q^* = \arg\min_{q \in Q \setminus C}\,
       r_z(C \cup \{q\})$.
\STATE Repeat Steps 2--3 until $|C| = k$.
\STATE \textbf{return} $C$
\end{algorithmic}
\end{algorithm}

\begin{algorithm}[t]
\caption{OC-Local$(X, Q, C, z)$}
\label{alg:oclocal}
\begin{algorithmic}[1]
\STATE For each center $c_i \in C$ and each unused pool point
       $q \in Q \setminus C$, compute $r_z(C - \{c_i\} + \{q\})$.
\STATE If any such swap strictly reduces $r_z(C)$, apply the first
       such swap: set $C \leftarrow C - \{c_i\} + \{q\}$ and
       restart Step 1.
\STATE If no improving swap exists, stop.
\STATE \textbf{return} $C$
\end{algorithmic}
\end{algorithm}

\section{The LP lower bound in full}
\label{app:lp}

For candidate radius $r$, the feasibility LP has variables
$y_i\in[0,1]$ (center mass at $i$), $s_j\in[0,1]$ (outlier mass
at $j$), and $x_{ij}\ge0$ for pairs with $d(x_i,x_j)\le r$:
\[
\begin{aligned}
  \textstyle\sum_{i:\,d(x_i,x_j)\le r}x_{ij}+s_j&\ge1 &&\forall j\\
  x_{ij}&\le y_i &&\forall i,j\\
  \textstyle\sum_i y_i&\le k\\
  \textstyle\sum_j s_j&\le z
\end{aligned}
\]
The smallest feasible $r$ satisfies $R_\mathrm{LP}\le\OPT$. We
binary-search the $O(N^2)$ pairwise distances. Solver: Gurobi~12.0.3,
feasibility tolerance $10^{-6}$. Adult LP times: $12$--$1{,}096$\,s at
$N\le2{,}000$; up to $3.4\times10^5$\,s at $N=20{,}000$.

\bibliographystyle{siamplain}
\bibliography{refs}

\end{document}
